\section{Related Work}
\label{sec:related}

\subsection{Evidence-based fact verification}
Automated fact-checking is commonly organized into claim detection, evidence retrieval, and claim verification; the last stage includes verdict prediction and, where applicable, justification production \citep{guo-etal-2022-survey}.
FEVER \citep{thorne-etal-2018-fever} makes evidence retrieval and verdict prediction measurable at scale.
Its 185{,}445 claims were generated by altering sentences extracted from Wikipedia, each claim is labeled Supported, Refuted or NotEnoughInfo, and for Supported and Refuted claims annotators recorded the Wikipedia sentences that serve as evidence.
The 2018 FEVER shared task, the first held on the dataset, asked systems to retrieve such evidence from Wikipedia and classify each claim \citep{thorne-etal-2018-fact}.
Its best system reached a FEVER score of 64.21\%, a score that credits a label only when the retrieved evidence is also correct, which shows that even with a Wikipedia corpus, finding the right evidence and reaching the right verdict remained difficult.
Later datasets extended claim verification beyond FEVER's setting: SciFact pairs expert-written scientific claims with research abstracts that support or refute them \citep{wadden-etal-2020-fact}; VitaminC builds contrastive claim--evidence pairs from Wikipedia revisions, in which nearly identical evidence supports a claim in one version and not in the other \citep{schuster-etal-2021-get}; HoVer requires evidence drawn from up to four Wikipedia articles \citep{jiang-etal-2020-hover}; and FEVEROUS annotates evidence as Wikipedia sentences, table cells, or both \citep{aly-etal-2021-feverous}.

Benchmark evidence does not guarantee that a model's verdict depends on it.
On FEVER, classifiers that see only the claim perform competitively with evidence-aware models by exploiting cues in the claim's wording \citep{schuster-etal-2019-towards}.
A model's accuracy when evidence is supplied therefore does not, by itself, show what the evidence contributed; that requires comparing the same claims with and without evidence.

\subsection{LLM reasoning with external evidence}
Supplying evidence to a large language model (LLM) does not ensure that the model uses it as intended.
In question answering, models can fall back on memorized knowledge when a supplied passage contradicts it \citep{longpre-etal-2021-entity}.
How well they use relevant information in a long input depends on where that information appears \citep{liu-etal-2024-lost}, and irrelevant material in the prompt can divert their reasoning \citep{shi2023distracted}.
When sources conflict, models weigh a document's relevance to the query heavily while largely ignoring stylistic features, such as scientific references or a neutral tone, that humans consider important \citep{wan-etal-2024-evidence}.
For fact verification on FEVER specifically, \citet{mamta-cocarascu-2025-facteval} find that LLM verdicts are brittle under small perturbations of the input.
Together, these studies show that evidence in the input can be ignored, outweighed or misread.
They do not measure how often adding correct, benchmark-designated evidence turns a correct answer into a wrong one.

\subsection{Evidence sufficiency and ambiguity}
Evidence that appears in a prompt is not necessarily evidence that a reader finds sufficient.
\citet{atanasova-etal-2022-fact} remove parts of the evidence supplied to fact-checking models and find that the models often fail to recognize when what remains is no longer enough to reach a verdict.
\citet{glockner-etal-2024-ambifc} document claim--evidence pairs that admit conflicting but valid interpretations and argue against forcing each claim into a single label.
For a study that asks readers to judge designated evidence, this means that each verdict must be interpreted within its limits.
A reader who reaches no decisive verdict has not shown that the evidence is insufficient, and a reader whose decisive verdict opposes the benchmark label has disagreed with that label, not shown it to be wrong.

\subsection{LLMs as judges}
An LLM judge is a language model prompted to evaluate an input against a rubric in place of a human rater.
LLM judges make large-scale evaluation feasible but bring their own biases.
\citet{zheng2023judging} document position, verbosity and self-enhancement biases; \citet{liu-etal-2023-g} raise the concern that LLM evaluators favor LLM-generated text; and \citet{chen-etal-2024-humans} find both human and LLM judges vulnerable to misinformation-oversight, gender, authority and beauty biases.
Position bias varies across judges and tasks rather than acting as a uniform artifact \citep{shi-etal-2025-judging}, and fine-tuned judge models are not a general substitute for GPT-4 as a judge \citep{huang-etal-2025-empirical}.
A description produced by one judge model therefore cannot be assumed to transfer to a judge from another model family.

\subsection{Gap addressed by this study}
Aggregate accuracy gains and existing research on evidence use leave two questions open.
The first is how often supplying correct, benchmark-designated evidence changes a correct answer into a wrong one.
Condition-level accuracy cannot show this, because gains on some claims offset losses on others.
The second is whether an automated description of those changes is stable when a judge from a different model family applies the same rubric.

Our experiment addresses the first question by supplying FEVER-designated gold evidence, the annotated evidence placed directly in the model input, and comparing each claim with and without it.
Because no retrieval step is involved, retrieval failure cannot explain a regression within this paired experiment.
Gold evidence does not, however, guarantee that the model attends to it, that every reader finds it sufficient, or that the internal cause of a regression is known.
The cost of the design is scope: it studies 2017-era Wikipedia claims restricted to FEVER's Supported and Refuted labels, with evidence that is designated rather than retrieved.

We address the second question with two judge panels, one from the Grok family and one from the Claude family.
Each panel consists of five isolated runs of a locked model that apply the same rubric to the same evidence; they are model-based judges, not independent human experts.
Where the two panels agree, a description of the regressions does not depend on one model family.
Where they disagree, the disagreement measures sensitivity to the judge family, and we report both panels rather than treating their two values as bounds on a true prevalence.
A panel's failure to reach a decisive verdict after all disclosure stages is recorded as an outcome in its own right, a model-judge analogue of the interpretive ambiguity documented by \citet{glockner-etal-2024-ambifc}.
